\SourceLeaf{130}
Wie sich übrigens jede Zahl $A$ unter obige Form bringen lasse, wird man leicht übersehen; $\mu$ und $a$ ergeben sich sogleich von selbst; macht man dann $A = 10^{\mu} a B$: so ist $b$ die Ziffer in der ersten Decimalstelle von $B$; setzt man ferner $B = \left(1+\frac{1}{10} b\right) C$, so ist $c$ die Ziffer in der zweiten Decimalstelle von $C$ u.\,s.\,w. Bei näherer Untersuchung findet sich, dass man noch bequemer $\frac{1}{A}$, ohne diesen Quotienten wirklich zu berechnen, unter die obige Form bringen kann, indem man successive $10^{\mu} a A = B$, $\left(1+\frac{1}{10} b\right) B = C$ u.\,s.\,w. macht, wo man also $B$ aus $A$, $C$ aus $B$ u.\,s.\,w. durch Multiplication erhält, indem bei der ersten Methode Division nöthig ist. \textsc{Leonelli} hat indessen den Calcul bei beiden Methoden durch Reduction auf einen bestimmten Mechanismus noch abgekürzt; wer häufigen Gebrauch von diesen Tafeln zu machen denkt, thut wohl, sich mit demselben vertraut zu machen. Nach einiger Übung wird man es dann gewiss bald dahin bringen, sich die anzuwendende Aufmerksamkeit mechanisch zu machen, und wir halten daher mit \textsc{Leonelli} den Vorwurf \textsc{Delambres} für übereilt, dass eine ermüdende Aufmerksamkeit nöthig sei, wenn man nicht jedesmal alle die vielen überflüssigen Nullen aufschreibe.

Man sieht leicht, dass man auf ähnliche Weise auch $A$ oder $\frac{1}{A}$ unter die Form
\[
10^{\mu} a \left(1+\frac{1}{100} b\right)\left(1+\frac{1}{10000} c\right)\ldots
\]
setzen könne, so dass $a$, $b$, $c\ldots$ lauter ganze Zahlen unter 100 bedeuten. Bei jedem einzelnen Falle wird man dann nur halb so viele Glieder nöthig haben; dagegen wird die Hülfstafel zwar auch nur halb so viel Ordnungen, aber in jeder 99 Logarithmen enthalten, also einen fast fünfmal so grossen Umfang haben müssen. \textsc{Leonelli} hat auch eine solche Hülfstafel, aber nur für die Briggischen Logarithmen, und nur auf 15 Decimalen beigefügt; sie enthält 485 Logarithmen. Es ist schade, dass nicht auch diese Tafel für 20 Decimalen eingerichtet ist.

Für die umgekehrte Aufgabe, zu einem gegebenen Logarithmen die Zahl zu finden, sind \textsc{Leonellis} Tafeln nicht weniger anwendbar. Durch Zerlegung des gegebenen Logarithmen in seine in den Tafeln befindlichen Bestandtheile erhält man die Factorenreihe der Zahl, und durch deren Multiplication (wofür sich leicht ein geschmeidiger Algorithmus findet) die Zahl selbst.

